% GPT-6 Astra Pro; /root/fixed_surface; 22 November 2026 UTC.
% Effort not exposed. New complete proof for independent review.
\section[HD-F1-S29: uniform obstruction on the geometric norm-source open]{HD-F1-S29: the actual norm sheaf fails weak semiregularity at every retained geometric fibre, without a period-rank restriction}
\label{sec:f1-norm-fourfold-uniform-obstruction}

Let $\mathcal U_{\rm geom}$ be the same nonempty open of the
fine cyclic-cover space fixed in S16: the base curve belongs
to the S13 canonical-divisor open, and the covering curve has
no $g^2_8$. These are the actual smooth-source and pencil
conditions used in S13--S22. Retain this whole open, before
the additional S23 period-rank restriction. In particular
$C$ is nonhyperelliptic. No nontrigonal condition is imposed.
After the finite etale component labeling already used there,
choose any of its three actual norm components and write
\[
 \pi:Z\to C,\qquad a:T\longrightarrow Y\hookrightarrow P,
 \qquad E=\mathcal O_Y=Ra_*\mathcal O_T.
\]
All conclusions below hold at every point of this specified
open. Put
\[
 \begin{gathered}
 U_C=H^0(C,K_C\varepsilon),\qquad
 V_C=H^0(C,K_C\varepsilon^{-1}),\\
 \mu_C:U_C\otimes V_C\longrightarrow H^0(C,K_C^2),
 \qquad \mu_C(s\otimes t)=st,\qquad
 r_C=\operatorname{rank}\mu_C.
 \end{gathered}
 \tag{F1S29.1}
\]
Here $\varepsilon$ is the original nontrivial order-three
line of the connected cover. The exact derivative $\eta_Y$
of S24 satisfies
\[
 \ker(\eta_Y|_{\mathfrak D_P})
       =\operatorname{im}d\operatorname{Pr},\qquad
 \operatorname{rank}(\eta_Y|_{\mathfrak D_P})=16-r_C\ge4.
 \tag{F1S29.2}
\]
Moreover the actual characteristic action induces an injection
\[
 \Phi_C:(\ker\mu_C)^*\lhook\joinrel\longrightarrow
  \ker\bigl(\sigma_E^2:\operatorname{Ext}^2_P(E,E)
                                      \longrightarrow HH_{-2}(P)\bigr).
 \tag{F1S29.3}
\]
Its image has dimension $16-r_C$. It is invariant for every
specified coherent finite point-action lift of $E$.
Thus this actual sheaf fails weak point-equivariant
semiregularity throughout $\mathcal U_{\rm geom}$, including
the trivial action. By S28 every faithful such finite point
action is in fact trivial. No assertion about the dimension
of the entire trace kernel is made.

\paragraph{1. The exact period sequence at every retained fibre.}
Riemann--Roch and nontriviality of $\varepsilon$ give
$\dim U_C=\dim V_C=4$, while $h^0(K_C^2)=12$.
The actual projector $1-(1+\sigma+\sigma^2)/3$ identifies
$H^1(P)$ with the nontrivial-character part of $H^1(Z)$.
The tangent space $\mathfrak D_P\subset H^1(P,T_P)$
preserving the specified polarization and order-three action
has dimension sixteen: polarization makes its two
opposite-character blocks transposes of one another, so one
arbitrary four-by-four block determines the tangent.
Additional endomorphisms at a special fibre do not enter
this deformation problem, which retains precisely the given
polarization and cyclic action.

As in S23, the actual Serre-dual identification is
\[
 \jmath:\mathfrak D_P\xrightarrow{\sim}(U_C\otimes V_C)^*,
 \qquad \jmath(\kappa)(s\otimes t)
                         =\langle\kappa(s),t\rangle_Z.
 \tag{F1S29.4}
\]
The tangent space of the fine cover space is $H^1(C,T_C)$,
by finite etale lifting across nilpotent thickenings.
This assertion requires no immersion property of its period
morphism. Variation of pulled-back forms and the actual trace
of the degree-three map give
\[
 \begin{split}
 \jmath(d\operatorname{Pr}(v))(s\otimes t)
   &=\langle\pi^*v,st\rangle_Z\\
   &=\langle v,\operatorname{Tr}_{\pi}(st)\rangle_C
     =3\langle v,\mu_C(s\otimes t)\rangle_C.
 \end{split}
 \tag{F1S29.5}
\]
Thus $\jmath\,d\operatorname{Pr}=3\mu_C^*$ under Serre
duality. This is also the multiplication formula of
Lange--Ortega, Proposition4.1; their subsequent injectivity
results are not needed. Dualizing $\mu_C$ gives the exact
sequence
\[
 0\longrightarrow(\operatorname{coker}\mu_C)^*
 \longrightarrow H^1(C,T_C)
 \xrightarrow{d\operatorname{Pr}}\mathfrak D_P
 \xrightarrow{q_C}(\ker\mu_C)^*\longrightarrow0,
 \qquad q_C(\kappa)=\jmath(\kappa)|_{\ker\mu_C}.
 \tag{F1S29.6}
\]
The factor three changes neither image nor kernel.
The period image has dimension $r_C$; the final quotient
has dimension $16-r_C$. This retains the first kernel
when $\mu_C$ is not surjective.

\paragraph{2. Uniform normal and characteristic maps.}
On all of $\mathcal U_{\rm geom}$, S16--S17 give the same
normal Cohen--Macaulay image and three possible pencil
local models. S21 gives the actual isomorphisms
$H^1(\mathcal O_P)\to H^1(\mathcal O_Y)$ and
$\operatorname{Alb}(T)\to P$. Their proofs do not use
the rank of $\mu_C$. The local resolutions and Postnikov
sequence of S24 therefore give the injections
\[
 H^1(Y,\mathscr N)\xrightarrow{j_N}
 \operatorname{Ext}^1_P(I_Y,E)
   \xrightarrow{d^{\rm ideal}_2}\operatorname{Ext}^2_P(E,E),
 \qquad d^{\rm ideal}_2(z)=z[1]\partial,
 \tag{F1S29.7}
\]
where $\partial:E\to I_Y[1]$ is the ideal-triangle boundary.
The second injection follows from the scalar retract
given by multiplication and the unit of $\mathcal O_Y$;
it needs no numerical scalar-cohomology assertion.

The map $\eta_Y$ is restriction followed by $H^1(\nu)$,
where
\[
 \nu:T_P|_Y\longrightarrow\mathscr N,\qquad
                         \nu(v)(f)=v(f)|_Y.
 \tag{F1S29.8}
\]
With diagonal convention $1\otimes f-f\otimes1\mapsto df$,
the first-jet lift of the unit gives
$\operatorname{At}(E)=(d|_{I_Y})[1]\partial$.
Consequently
\[
 \operatorname{ob}_E(\kappa)
  =(1_E\otimes\kappa)\operatorname{At}(E)
       =d^{\rm ideal}_2j_N\eta_Y(\kappa).
 \tag{F1S29.9}
\]
S25's deformation constructions apply on the same geometric
open: zero $\eta_Y(\kappa)$ is equivalent to a flat embedded
lift in the specified first-order ambient deformation, and
such a lift has the smooth canonical resolution proved there.
These statements precede the rank-twelve specialization
of the period quotient.

\paragraph{3. Every actual cover tangent is in the normal kernel.}
Take $v\in H^1(C,T_C)$ and an actual marked curve deformation
over $D=\operatorname{Spec}\mathbf C[\epsilon]/(\epsilon^2)$.
The original finite etale cover extends uniquely, with its
deck map and group law. This deformation remains in
$\mathcal U_{\rm geom}$ because that locus is open.

The source family can be constructed before this first-order
pullback. On a smooth local chart of the fine cover space,
the total relative symmetric products are smooth of equal
dimension. Their finite map is flat by the finite
Cohen--Macaulay flatness criterion. Base change to the
relative canonical $\mathbf P^4$-bundle makes the relative
norm fibre flat. All its geometric fibres are smooth by S13,
so the family is proper smooth. Its connected components
are labeled by an etale base change. The chosen component
therefore pulls back to an actual smooth proper $T_D$.
An origin in the smooth relative Prym torsor lifts, giving
the actual relative Abel morphism
\[
                          a_D:T_D\longrightarrow P_D.
 \tag{F1S29.10}
\]
No submersion property of the period map is used here.

The embedded quotient is the S25 unit construction.
Applying $Ra_{D*}$ to
$0\to\mathcal O_T\xrightarrow{\epsilon}\mathcal O_{T_D}
\to\mathcal O_T\to0$, using $Ra_*\mathcal O_T=E$, gives
$R^{>0}a_{D*}\mathcal O_{T_D}=0$ and
$0\to E\to F_D\to E\to0$, where
$F_D=a_{D*}\mathcal O_{T_D}$.
Multiplication by $\epsilon$ is the last surjection followed
by the first inclusion, so $\ker\epsilon=\epsilon F_D$ and
$F_D$ is $D$-flat. The algebra unit
$\mathcal O_{P_D}\to F_D$ is surjective by reduction and
Nakayama. It gives an actual embedded $Y_D\subset P_D$.
The gluing criterion in (F1S29.8) yields
\[
                       \eta_Y\,d\operatorname{Pr}(v)=0.
 \tag{F1S29.11}
\]

\paragraph{4. Every normal-kernel tangent recovers a cover tangent.}
Take $\kappa\in\ker(\eta_Y|_{\mathfrak D_P})$.
S25 gives $Y_D\subset P_D$ and its smooth canonical
resolution $T_D$. By S28, all of S26's marked-divisor
and relative Hilbert conclusions apply on this whole
geometric open. They give an actual smooth Hilbert
curve $Z_D$ and its universal divisor, finite flat of
degree eight over $T_D$.

Here are the S27 reconstruction maps, with their
rank-independent hypotheses retained. The pullback
$P_D^\vee\to\operatorname{Pic}^0(T_D/D)$ reduces to the
S19/S21 isomorphism and is an isomorphism by flatness and
Nakayama. Dualizing gives
$A_D:\operatorname{Alb}(T_D/D)\simeq P_D$.
The universal degree-eight divisor, subtracting its
value at a section $t_0$, gives $u_D:T_D\to J_D=J(Z_D/D)$
and its Albanese homomorphism $U_D$. Thus
\[
                  i_D=U_DA_D^{-1}:P_D\hookrightarrow J_D
 \tag{F1S29.12}
\]
is the closed immersion reducing to the literal original
inclusion $i:P\hookrightarrow JZ$.

For the principal Jacobian polarization $\lambda_D$ and
the given action $\omega_D$ on $P_D$, homomorphism rigidity
gives $i_D^\vee\lambda_Di_D=\beta_D$. Put
\[
 r_D=\beta_D^{-1}i_D^\vee\lambda_D,\qquad
 s_D^{\rm rat}=1+i_D(\omega_D-1)r_D.
 \tag{F1S29.13}
\]
The self-adjoint projector $i_Dr_D$ makes this action
polarization-preserving and of order three. Centrally
it is the original deck map $\sigma$: identity on
$\pi^*JC$ and $\omega$ on $P$.
The type $(1^4,3^4)$ makes $3\beta_D^{-1}$ integral.
Hence $3s_D^{\rm rat}$ is integral and is zero on
$J_D[3]$, since its central restriction is zero and
this group is finite etale. Factoring through the
actual quotient $[3]$ gives an integral $s_D$ with
this rationalization. The torsion-free Hom group
gives its integral order-three and polarization
identities, exactly as in S27.

The free order-three action makes $Z$ nonhyperelliptic:
a hyperelliptic quotient would have two invariant
fibres of degree two over the fixed points of the
induced order-three action on $\mathbf P^1$, contradicting
freeness. Infinitesimal Torelli for this $Z$ is injective;
its dual is Max Noether's surjective canonical
multiplication. The extension $s_D$ makes the Jacobian
deformation class $\sigma$-invariant. Functoriality
and injectivity therefore make the class $\xi_Z$
of $Z_D$ invariant. The actual etale differential
and averaging give
\[
 \pi^*:H^1(C,T_C)\xrightarrow{\sim}
                      H^1(Z,T_Z)^{\langle\sigma\rangle}.
 \tag{F1S29.14}
\]
Choose its inverse image $\xi_C$ and an actual $C_D$.
Finite etale lifting extends $\pi$ over $C_D$; its
upstairs class is $\pi^*\xi_C=\xi_Z$. The marked
first-order classification and $H^0(Z,T_Z)=0$ identify
this upstairs curve with $Z_D$, retaining its cyclic
action and group law. The induced Jacobian action
equals $s_D$ by rigidity.

For $q_D:J_D\to J_D/i_D(P_D)$, rigidity gives
$q_D(1-s_D)=0$. Conversely
$(1-s_D)i_D=i_D(1-\omega_D)$, and $1-\omega_D$
is an isogeny. Thus the actual image of $1-s_D$
is exactly $i_D(P_D)$, the connected norm kernel
of the recovered cover. Its polarization and action
are the original specified $\beta_D,\omega_D$.
Consequently its marked tangent is exactly $\kappa$:
\[
                            \kappa=d\operatorname{Pr}(\xi_C).
 \tag{F1S29.15}
\]
Together with (F1S29.11), this proves the kernel
equality in (F1S29.2) everywhere on the retained
geometric open. It never requires uniqueness of
a preimage under $d\operatorname{Pr}$.

\paragraph{5. The whole character and the exact trace-kernel injection.}
The entire S14 character holds on this geometric open:
\[
 \operatorname{ch}(E)=w+\frac{\beta^4}{72}
       -\frac{\beta^5}{120}+\frac{17\beta^6}{6480}
       -\frac{13\beta^7}{22680}+\frac{53\beta^8}{544320},
 \qquad w\in W_K(P).
 \tag{F1S29.16}
\]
S14 used the actual smooth proper source family,
relative GRR and proper cohomological Gysin to
make this a horizontal class, then the connected
monodromy calculation. It used no period-rank
hypothesis. This is an identity of flat cohomology
local systems on the same component and persists
at every retained point. Additional Hodge classes
on a special fibre do not alter this given class.

Choose a $\beta$-unitary eigenbasis
$a_1,\ldots,a_4,b_1,\ldots,b_4$ at the given point,
and put $A=a_1\cdots a_4$, $B=b_1\cdots b_4$.
Then $w=\lambda A\bar B+\bar\lambda B\bar A$.
The Weil tangents are spanned by
\[
 \kappa_{ij}=\partial_{a_i}\otimes\bar b_j
                       +\partial_{b_j}\otimes\bar a_i.
 \tag{F1S29.17}
\]
They annihilate $\beta$ by cancellation of the two
antiholomorphic terms. Contraction with $A\bar B$
repeats a factor of $\bar B$; contraction with
$B\bar A$ repeats a factor of $\bar A$.
They therefore annihilate $w$ and every power of
$\beta$. The full trace-characteristic identity
with the same diagonal Atiyah action gives
\[
 \sigma_E^2\operatorname{ob}_E(\kappa)
      =\kappa\mathbin{\lrcorner}\operatorname{ch}(E)=0
                         \quad(\kappa\in\mathfrak D_P).
 \tag{F1S29.18}
\]
The proved normal-kernel equality and (F1S29.6)
make the descended map
$\overline\eta_Y:(\ker\mu_C)^*\to H^1(Y,\mathscr N)$
injective. The promised map is
\[
 \Phi_C=d^{\rm ideal}_2j_N\overline\eta_Y,\qquad
           \Phi_Cq_C=\operatorname{ob}_E|_{\mathfrak D_P}.
 \tag{F1S29.19}
\]
All three factors are injective; (F1S29.18) supplies
the full trace-kernel target. Its image dimension
is $16-r_C\ge4$, proving (F1S29.2)--(F1S29.3).
For $r_C=12$ this recovers S27's four-dimensional
subspace. The proof also applies at any retained
point with smaller $r_C$, without asserting that
such a point exists in this particular open.

\paragraph{6. The specified finite action and exact source scope.}
For an actual point functor $F_g=t_x^*(-)\otimes L$,
translation acts identically on $H^1(P,T_P)$,
and $\operatorname{At}(L)=0$ for $L\in\operatorname{Pic}^0(P)$.
Atiyah-class naturality under the specified
isomorphism $F_gE\simeq E$ fixes every class
in (F1S29.19) under the actual conjugation action
on self-Ext. The identity-sector induction/forgetful
comparison of R05, with the same functors and lift,
carries this nonzero invariant kernel into the
weak $G$-semiregularity kernel. No coherence data
or lift are replaced.

The fresh period-source check is the complete
Section4 proof of Lange--Ortega, \emph{Prym varieties
of cyclic coverings}, Geom. Dedicata150 (2011),
391--403, DOI10.1007/s10711-010-9512-9, including
Proposition4.1 and Lemmas4.2,4.5. Its tangent
formula holds at every smooth cover in its
stated moduli problem. The infinitesimal Torelli,
Noether multiplication and finite-etale lifting
inputs are the exact complete primary proofs
specified and checked in S27. The S26--S28
Fourier and Hilbert input has the actual enlarged
nonhyperelliptic scope proved there. All original
smoothness and pencil conditions are retained;
the result makes no extension across the other
excluded geometric loci and no assertion about
another perfect object.

